% !TeX root = ../sdc_regulations.tex
% Add the above to each chapter to make compiling the PDF easier for some editors.

\newpage
{\bf Preface}
\\ \\ Welcome to the Swarm Drone Challenge! \\ \\
The Swarm Drone Challenge 2026 (SDC26) is an international competition dedicated to advancing the capabilities of autonomous aerial systems through the development and demonstration of swarm intelligence. It invites students, professionals, and startups to design and implement self-developed algorithms that enable groups of drones to cooperate effectively in complex, realistic mission environments.

Building on the success of previous editions, SDC26 continues the established Capture the Flag scenario, in which two teams deploy multiple drones operating collaboratively as a swarm. This competitive format promotes diverse technical and strategic approaches, addressing key challenges such as resilient communication, AI-based coordination, adaptive decision-making, and collective mission behavior under dynamic conditions.

The purpose of the Swarm Drone Challenge is to explore novel methods for distributed autonomy, evaluate innovative algorithmic concepts, and contribute to the evolution of common frameworks and standards for cooperative unmanned systems. By fostering experimentation and knowledge exchange among participants, SDC26 aims to identify promising technological directions and advance the understanding of swarm robotics in both theoretical and applied contexts.
\newpage

\newpage
{\bf Changelog}
\\ \\{\bf Changes in version: \getVersion{}}
\begin{itemize}
	\item{Defined tournament structure: 6 teams compete across three stages (Opening Round, Semi-Finals, Final); the 4 best-ranked teams advance from the Opening Round and the 2 match winners advance to the Final. Pairings are set by the organizers and announced via Discord.}
	\item{Defined overall challenge scoring: maximum 25 challenge points across three categories — up to 15 from match wins (5 per won match), up to 5 from jury evaluation (pitch and overall approach), and up to 5 from the technical report.}
	\item{Defined technical report requirements: must cover localisation approach, flight control and trajectory planning strategies, and swarm strategy. Report is due 4 days before the Final. Open-sourcing code is considered positively.}
	\item{Removed the rule that target object positions vary per match.}
	\item{Confirmed challenge location: ILA Berlin, Thursday 11 June 2026.}
	\item{Differentiated between game points (scored during a match via box captures: 1, 5, or 10) and challenge points (overall tournament leaderboard: match wins, jury evaluation, technical report).}
	\item{Fixed ArUco marker positions and IDs in the arena.}
	\item{Corrected field length from 25\,m to 20\,m.}
	\item{Removed zero ArUco marker from positioning system.}
	\item{Restructured document into four sections: Competition Rules, Team Regulations, Safety Regulations, and Technical Regulations.}
	\item{Promoted gameplay objective to introductory prose in the Target Objects subsection.}
	\item{Split Intellectual Property and Reporting Requirements into separate subsections.}
	\item{Consolidated all safety-related rules into a dedicated Safety Regulations section.}
	\item{Merged Drone Size/Weight Limits and Drone Identification into a single Drone Requirements subsection.}
	\item{Added Game Logic section with precise capture mechanics, ArUco marker encoding table, point scoring rules (1/5/10 points), and Special Maneuver definition.}
\end{itemize}
\newpage
